\documentstyle[%


righttag%


,11pt%


,amssymb%


,russian%               remove in Eng.


,izvru%                 Set izven% in Eng.


]{amsart}





%





\newcommand{\tts}[1]{}





\theoremstyle{plain}


\theorembodyfont{\it}


\newtheorem{thms}{\hskip\parindent Теорема}[section]


\newtheorem{lems}{\hskip\parindent Лемма}[section]


\newtheorem{props}{\hskip\parindent Предложение}[section]


\newtheorem{probs}{\hskip\parindent Задача}[section]





\theoremstyle{definition}


\theorembodyfont{\rm}


\newtheorem{cors}{\hskip\parindent Следствие}[section]





%\hoffset=-15mm%                  For Eng.


\setcounter{page}{71}


\god{2003}


\nomer{12~(499)} %                        Remove in Eng.


%\nomer{Vol.\,47, No.\,12}%               Set in Eng.


%\str{\pageref{begin}--\pageref{end}}%  Set in Eng.


\udk{517.853}





\author{\it И.В.\,Коннов, О.В.\,Пинягина}


% NB:   ^^ remove in Eng.


\title{Метод спуска по интервальной функции\\ для негладких задач равновесия}


\thanks{Работа выполнена при финансовой поддержке Российского фонда 


фундаментальных исследований (коды проектов


\No\No\,01-01-00068, 03-01-96012), фонда НИОКР Академии наук 


Республики Татарстан и Академии наук Финляндии (проект \No\,77796).}





\date{}


%\pagestyle{myheadings}%         Set in Eng.


\pagestyle{plain} %              Remove in Eng.


\begin{document}


%\thispagestyle{plain}%          Set in Eng.


\maketitle


\label{begin}











\section{Введение}





Пусть $U$ --- непустое замкнутое выпуклое множество в пространстве $R^n$,


$\Phi : R^n \times R^n \rightarrow R$ --- равновесная функция,


т.\,е. $\Phi(u,u)=0$ для любого $u \in U$. 


Задача равновесия определяется следующим образом: 


найти точку $u^* \in U$ такую, что


$$


\Phi(u^*,v) \ge 0 \quad \forall v \in U.


$$


Задачи равновесия позволяют единым образом формулировать 


и исследовать разнообразные сложные проблемы, возникающие 


в экономике, математической физике, исследовании операций и других областях.


Более того, они тесно связаны с другими общими задачами нелинейного анализа, 


как, например, задачи оптимизации, задачи о дополнительности и 


вариационные неравенства. По этой причине теория и методы решения 


задач равновесия


довольно интенсивно исследуются \cite{BC84}--\cite{KS00}.


Как известно, один из наиболее распространенных подходов к 


решению вариационных


неравенств состоит в сведении их к соответствующей оптимизационной задаче 


с помощью так называемых интервальных (или оценочных) функций


\cite{Pat97}. Хотя эти функции, как правило, являются невыпуклыми, они 


позволяют преодолеть многие трудности как в теории, так и при построении 


методов решения вариационных неравенств. 





Главная цель данной работы состоит в применении аппарата интервальных функций 


к довольно широкому классу задач равновесия, включающих негладкие функции. 


На этой основе исходная задача заменяется эквивалентной задачей о необходимых 


условиях оптимальности в негладкой оптимизации.


Кроме того, показано, что можно построить итерационный метод спуска 


по интервальной функции, сходящийся к решению задачи при некоторых 


дополнительных предположениях.








\section{Интервальная функция для задач равновесия}





Пусть $f : R^n \rightarrow R$ --- выпуклая, но не обязательно 


дифференцируемая функция, 


$h :\linebreak R^n \times R^n \rightarrow R$ --- равновесная функция. 


Будем также предполагать, что $h$ --- 


дифференцируемая функция и $h(u,\cdot)$ --- выпуклая функция 


для любого $u \in R^n$.





Рассмотрим задачу равновесия в следующей форме \cite{BC84}. 





\begin{probs}\label{n1}


Найти точку $u^* \in U$ такую, что


$$


h(u^*,v) + f(v) - f(u^*) \ge 0 \quad \forall v \in U.


$$


\end{probs}





\noindent{}Множество решений задачи \ref{n1} будем обозначать через $U^*$, 


а через 


$\nabla_u h(\cdot, \cdot)$ ($\nabla_v h(\cdot, \cdot)$) ---


частный градиент функции $h$ по первому (второму) векторному аргументу.





Вначале сформулируем условие оптимальности для задачи \ref{n1} 


в форме смешанного вариационного неравенства.





\begin{props}\label{p11}


Точка $u^* \in U$ является решением задачи $\ref{n1}$


тогда и только тогда, когда 


\begin{equation} \label{n1a}


\langle \nabla_v h(u^*,u^*), v - u^*\rangle +


f(v) - f(u^*) \ge 0 \quad \forall v \in U.


\end{equation}


\end{props}





\begin{pf}


{\bf Необходимость.} Так как $u^* \in U^*$, то $u^*$ 


является решением задачи выпуклого программирования


$$


\min \limits_{v \in U} \rightarrow h(u^*,v) + f(v).


$$


Записав условия оптимальности для этой задачи, получим


$$


\langle \nabla_v h(u^*,u^*) + g, v - u^* \rangle\ge0 \quad \forall v \in U,


$$


где $g$ --- субградиент функции $f$  в точке $u^*$. По определению 


субградиента


$f(v) - f(u^*) \ge\linebreak \langle g, v - u^* \rangle$. Поэтому для $u^*$ 


выполняется условие (\ref{n1a}).


\medskip





{\bf Достаточность.} Предположим, что для $u^*$ выполняется условие 


(\ref{n1a}). Принимая во внимание, что $h(u^*,\cdot)$ выпукла, получим 


$$


\langle \nabla_v h(u^*,u^*), v - u^* \rangle


\le h(u^*,v) - h(u^*, u^*) = h(u^*,v)


\quad \forall v \in U.


$$


Отсюда в силу условия (\ref{n1a}) имеем


$h(u^*,v) + f(v) - f(u^*) \ge 0$ $\forall v \in U$,


т.\,е. $u^* \in U^*.$


\end{pf}








Рассмотрим модифицированную по отношению к \ref{n1} задачу равновесия. 





\begin{probs}\label{n2}


Найти точку $\ov{u} \in U$ такую, что 


$$


h(\ov{u},v) + f(v) - f(\ov{u}) + 0.5 \alpha \|v-\ov{u}\|^2 \ge 0


\quad \forall v \in U,


$$


где $\alpha > 0$ --- заданное число.


\end{probs}





Обозначим через $U^\alpha$ множество решений задачи \ref{n2}.





\begin{props}\label{p12} 


Множества решений задач $\ref{n1}$ и $\ref{n2}$ совпадают,


т.\,е. $U^* = U^\alpha$.


\end{props}





\begin{pf}


Результат следует непосредственно из предложения \ref{p11}, 


поскольку \linebreak


$\nabla_v h(u,u) = \nabla_v \ov h(u,u)$,


где $\ov h(u,v) = h(u,v) + 0.5 \alpha \|v-u\|^2$.


\end{pf}





Введем обозначения


\begin{gather}


L(u,v) = - h(u,v) - f(v) + f(u) - 0.5 \alpha \|v-u\|^2, \notag \\


\label{n3}


\psi(u)= \sup \limits_{v \in U} L(u,v) = L(u, v(u)). 


\end{gather}


Заметим, что точка $v(u)$ всегда существует и является единственным решением


задачи (\ref{n3}), поскольку $L(u,\cdot)$


--- сильно вогнутая и непрерывная функция. Условие оптимальности для этой 


задачи можно сформулировать следующим образом.





\begin{props}\label{p13}


Для всех $u' \in U$ верно неравенство


\begin{equation}\label{n3a}


\langle \nabla_v h(u',v(u')) + \alpha (v(u') - u'), v - v(u') \rangle +


f(v) - f(v(u')) \ge 0 \quad \forall v \in U.


\end{equation}


\end{props}





\begin{pf}


Выберем любую точку $v \in U$  и 


для краткости обозначим $v' = v(u')$. 


Записав условия оптимальности для задачи (\ref{n3}), имеем


$$


\langle -\nabla_v h(u',v') - \alpha (v' - u') - g, v' - v\rangle  \ge 0, 


$$


где $g$ --- субградиент функции $f$ в точке $v'$. Отсюда согласно определению 


субградиента получим


$$


\langle \nabla_v h(u',v') + \alpha (v' - u'), v - v'\rangle  


+ f(v) - f(v') \ge 0,


$$


т.\,е. (\ref{n3a}) выполняется.


\end{pf}





Покажем, что определенная в (\ref{n3}) функция $\psi$ является


интервальной в задаче равновесия \ref{n1},


т.\,е. что имеет место





\begin{props}\label{p14}


Функция $\psi$ обладает следующими свойствами\/${:}$


\begin{itemize}


\item[(i)] $\psi(u) \ge 0$ $\forall u \in U;$


\item[(ii)] $\psi(u) = 0$ и $u \in U \Longleftrightarrow u \in U^*$.


\end{itemize}


\end{props}





\begin{pf}


Утверждение (i) следует из того факта, что $L(u,u)=0$.


Докажем утверждение (ii). Предположим, что $u \in U$ и $\psi(u) = 0$, тогда 


$$


-h(u,v) - f(v) + f(u) - 0.5 \alpha \|v-u\|^2 \le 0 \quad 


\forall v \in U,


$$


т.\,е. $u \in U^\alpha$ и $u \in U^*$ в силу предложений \ref{p11} 


и \ref{p12}.





Предположим теперь, что $u \in U^*$. Тогда по определению


$$


-h(u,v) - f(v) + f(u) \le 0 \quad \forall v \in U.


$$


Следовательно, 


$$


-h(u,v) - f(v) + f(u) - 0.5 \alpha \|v-u\|^2 \le 0 


\quad \forall v \in U.


$$


С другой стороны, в силу утверждения (i) 


$$


\psi(u) \ge 0 \quad \forall u \in U.


$$


Отсюда следует $\psi(u) = 0$.


\end{pf}





Приведем еще одну формулировку условия оптимальности для задачи 


\ref{n1} в форме неподвижной точки отображения $u \mapsto v(u)$.





\begin{props}\label{p15}


$u^* =v(u^*) \Longleftrightarrow u^* \in U^* $.


\end{props}





\begin{pf}


Предположим, что $u^* \in U^*$, тогда вследствие предложения \ref{p12} имеем 


$u^* \in U^\alpha$ и в силу предложения \ref{p11} выполняется


$$


\langle \nabla_v h(u^*,u^*), v(u^*) - u^*\rangle  + f(v(u^*)) - f(u^*) \ge 0.


$$


В то же время в силу предложения \ref{p13}


$$


\langle \nabla_v h(u^*,v(u^*)), u^* - v(u^*)\rangle  + f(u^*) - f(v(u^*)) - 


\alpha\|u^* - v(u^*)\|^2 \ge 0.


$$


Сложив оба неравенства, получим


$$


-\alpha\|u^* - v(u^*)\|^2 \ge \langle \nabla_v h(u^*,v(u^*)) -


\nabla_v h(u^*,u^*), v(u^*) - u^*\rangle  \ge 0,


$$


поскольку $\nabla_v h(u^*,\cdot)$ является монотонным отображением. 


Отсюда следует $u^* = v(u^*)$. Теперь докажем обратное. 


Предположим, что $u^* = v(u^*)$, тогда по определению $u^* \in U$ и


$$


\psi(u^*) = L(u^*, u^*) = 


- h(u^*,u^*) - f(u^*) + f(u^*) - 0.5\alpha\|u^* - u^*\| = 0,


$$


отсюда в силу (ii) предложения \ref{p14} \ $u^* \in U^*$. 


\end{pf}





Свойства функции $\psi$, доказанные в предложении \ref{p14}, 


т.\,е. неотрицательность на допустимом множестве и равенство нулю на 


множестве решений, показывают, что $\psi$ может служить интервальной функцией 


для задачи \ref{n1}. Отсюда, в частности, следует, что исходная задача 


\ref{n1} эквивалентна задаче условной минимизации 


\begin{equation}\label{minpsi}


\min \limits_{u \in U} \longrightarrow \psi(u).


\end{equation}


Однако выпуклость функции $\psi$ не гарантируется и данная задача может иметь


в принципе локальные минимумы, отличные от глобальных. Поэтому было бы удобнее


заменить эту задачу на ее условия стационарности. Это потребует дополнительных 


предположений, которые изучаются в следующем параграфе.





\section{Непрерывность и стационарность}





В дальнейшем будут использованы следующие предположения.


\begin{itemize}


\item[(A1)] Для отображения $\nabla_v h(\cdot,v)$ выполняется условие Липшица 


с константой $L_h$ для любого $v \in R^n$.


\item[(A2)] Для любых $u, v \in R^n$ выполняется неравенство


$$


\langle \nabla_u h(u,v) + \nabla_v h(u,v), v - u\rangle


\ge \varkappa \|u - v\|^2,


$$


где $\varkappa > 0$.


\end{itemize}





Отметим, что в случае вариационного неравенства, т.\,е. когда функция $h$ 


представлена в виде $h(u,v)= \langle G(u), v-u \rangle$,


свойство (А2) выглядит следующим образом:


$$


\langle \nabla G(u)^T(v-u),v-u \rangle \ge \varkappa \|u-v\|^2,


$$


что соответствует требованию сильной монотонности отображения $G$.





\begin{lems} \label{l21}


Если условие {\em (A1)} выполняется, то отображение $u \mapsto v(u)$ 


удовлетворяет условию Липшица.


\end{lems}





{\bf Доказательство.}


Выберем произвольные точки $u', u'' \in U$ 


и обозначим $v' = v(u')$, $v'' = v(u'')$. 


Тогда в силу предложения \ref{p13} имеем 


\begin{align*}


\langle \nabla_v h(u',v') + \alpha(v' - u'), v'' - v'\rangle  + 


f(v'') - f(v') &\ge 0, \\


%


\langle \nabla_v h(u'',v'') + \alpha(v'' - u''), v' - v''\rangle  + 


f(v') - f(v'') &\ge 0.


\end{align*}


Сложив эти неравенства, получим


$$


\langle \nabla_v h(u',v') - \nabla_v h(u'',v''), v'' - v'\rangle  + 


\alpha \langle (v' - v'') + (u'' - u'), v'' - v'\rangle  \ge 0.


$$


Отсюда следует


\begin{align*}


\langle \nabla_v h(u',v') - \nabla_v h(u'',v'), v'' - v'\rangle  + 


\alpha \langle u'' - u', v'' - v'\rangle  & \ge \\


\ge \langle \nabla_v h(u'',v') - \nabla_v h(u'',v''), v' - v''\rangle  + 


\alpha \|v'' - v'\|^2 & \ge \alpha\|v'' - v'\|^2,


\end{align*}


поскольку $\nabla h_v(u'',\cdot)$ монотонно. Далее, используя (A1), получим


$$


(L_h + \alpha) \|u' - u''\|\, \|v' - v''\| \ge \alpha\|v' - v''\|^2,


$$


отсюда


$$


\|v' - v''\| \le (1 +L_h/\alpha) \|u' - u''\|.\quad\square


$$





Из леммы \ref{l21} и формулы (\ref{n3}), в частности, следует, что функция


$\psi$ непрерывна на $U$.





\begin{lems}\label{l22} 


Пусть выполняется условие {\em (A1)}. Тогда функция $\psi$ имеет производную


в любой точке $u \in U$ по любому направлению $d \in R^n$, причем


$$


\psi'(u;d) = f'(u;d) - \langle \nabla_u h(u, v(u)) + \alpha (u-v(u)),d \rangle.


$$


\end{lems}





{\bf Доказательство.} 


Выберем произвольно вектор $d \in R^n$. 


Учитывая свойство непрерывности отображения $u \mapsto v(u)$ и определение 


функции $\psi$, получаем, что в данных условиях можно применить теорему 3.4


из \cite{DR90} о дифференцировании функции максимума, откуда следует, что


функция $\psi$ дифференцируема по направлениям в любой точке $u \in U$ и 


$$


\psi'(u;d) = f'(u;d) - \langle \nabla_u h(u, v(u)) +


\alpha (u-v(u)),d \rangle.\quad\square


$$





Сформулируем необходимое и достаточное условие оптимальности 


для задачи (\ref{minpsi}).





\begin{thms}[{\series{m}\shape{n}\selectfont условие стационарности}]


\label{tst}


Пусть выполняются условия {\em (A1)} и {\em (A2)}. Тогда 


$$


\psi'(u';u-u') \ge 0 \quad \forall u \in U \Longleftrightarrow u' \in U^*.


$$


\end{thms}





\begin{pf}


То, что решение задачи \ref{n1} должно


удовлетворять условию неотрицательности производной по любому направлению 


в этой точке, является очевидным. 


Предположим теперь, что 


$$


\psi'(u';u-u') \ge 0 \quad \forall u \in U.


$$ 


Из определения производной по направлению следует


$$


f'(u';u-u') - \langle \nabla_u h(u', v(u')) + 


\alpha (u'-v(u')),u-u' \rangle \ge 0 \quad \forall u \in U.


$$


Полагая здесь $u=v(u')$, получим


$$


f'(u';v(u')-u') -  \langle \nabla_u h(u', v(u')) + 


\alpha (u'-v(u')),v(u')-u' \rangle \ge 0.


$$


С другой стороны, в силу предложения \ref{p13}, полагая в (\ref{n3a}) 


$v=u'$, имеем


$$


\langle \nabla_v h(u',v(u')) + \alpha (v(u') - u'), u' - v(u') \rangle  + 


f(u') - f(v(u'))  \ge 0.


$$


Складывая два последних неравенства, получим


\begin{equation} \label{st1}


f'(u'; v(u')-u') +f(u') - f(v(u')) - 


\langle \nabla_u h(u',v(u')) + \nabla_v h(u',v(u')), v(u') -


u' \rangle \ge 0.  


\end{equation}


Сумма первых трех слагаемых неположительна вследствие выпуклости функции $f$,


поэтому в силу предположения (A2) из (\ref{st1}) следует 


$ -\varkappa \|u'-v(u')\| \ge 0,$ что выполняется только если 


$u'=v(u')$. Тогда согласно предложению \ref{p15} следует $u' \in U^*.$ 


\end{pf}





Итак, любая стационарная точка задачи (\ref{minpsi}) 


дает решение исходной задачи 


равновесия \ref{n1}, поэтому для отыскания решения можно использовать 


соответствующие итеративные методы минимизации с учетом негладкости 


функции $\psi$. Отметим, что из теоремы \ref{tst} 


и предложения \ref{p14} следует, что условие 


(A2) фактически гарантирует отсутствие локальных 


минимумов функции $\psi$, отличных от глобального.





В следующем разделе покажем, что если точка $u$ не является решением 


задачи \ref{n1}, то вектор $d = v(u) - u$ представляет собой направление 


убывания для функции $\psi$ в точке $u$. 


Основываясь на этом утверждении, построим


метод спуска для решения задачи \ref{n1}.





\section{Метод спуска}





\begin{lems}\label{ldesc}


Пусть выполняются условия {\em (A1)} и {\em (A2)}. Тогда для любого 


$u \in U \setminus U^*$


$$


\psi'(u; v(u)-u) < 0.


$$


\end{lems}





\begin{pf}


Выберем любую точку $u \in U$. Согласно лемме \ref{l22}


$$


\psi'(u;v(u)-u) = 


f'(u;v(u)-u) - \langle \nabla_u h(u, v(u)) + 


\alpha (u-v(u)),v(u)-u \rangle.


$$


С другой стороны, в силу предложения \ref{p13}, полагая в (\ref{n3a})


$u=u'$, $v=u$, имеем


$$


0 \le \langle \nabla_v h(u,v(u)) + \alpha (v(u) - u), u - v(u) \rangle  + 


f(u) - f(v(u)).


$$


Складывая эти соотношения, получим


$$


\psi'(u;v(u)-u) \le f'(u; v(u)-u) +f(u) - f(v(u)) -


\langle \nabla_u h(u,v(u)) + \nabla_v h(u,v(u)), v(u) - u \rangle.  


$$


Для любой точки $u$ из множества $U \setminus U^*$ сумма первых трех 


слагаемых в правой части неравенства неположительна 


вследствие выпуклости функции $f$, четвертое слагаемое отрицательно


в силу предположения (A2), поскольку $u \not= v(u)$. Таким образом, 


$\psi'(u; v(u)-u) < 0$.


\end{pf}





Итак, решение задачи (\ref{n3}) позволяет получить направление убывания 


функции $\psi$ в любой точке $u \in U$.





Построим теперь для решения задачи \ref{n1}


\medskip





{\bf Метод спуска}





{\it Шаг} 0. 


Выберем произвольно точку $u^0 \in U$ и параметр $\alpha >0$. Положим $k=0$.





{\it Шаг} 1. 


Вычислим $\psi(u^k)$ и $v(u^k)$ по формуле 


$$


\psi(u^k)= \sup \limits_{v \in U} L(u^k,v) = L(u^k, v(u^k)).


$$ 


Если $\psi(u^k)=0$, то $u^k$ является решением 


задачи \ref{n1}, процесс решения останавливается.





{\it Шаг} 2. 


Вычислим $w(u^{k})=u^k+l_kd^k$, где $d^k=v(u^k)-u^k,$ а


$l_k$ представляет собой решение задачи одномерной минимизации


$$


\min \limits_{0 \le l \le 1} 


\longrightarrow \psi(u^k+ld^k),


$$


положим $u^{k+1}=w(u^k)$, заменим $k$ на $k+1$ и перейдем к шагу 1.


\medskip





Введем дополнительное предположение.


\begin{itemize}


\item[(A3)] Для любого $v \in R^n$ отображение 


$\nabla_v h(\cdot, v)$ является сильно монотонным на множестве $U$


с константой $\tau.$


\end{itemize}





\begin{thms}\label{terr}


Пусть выполняется предположение {\em (A3)}. Тогда найдется 


число $\sigma >0$ такое, что


$$


\sigma \|u-u^*\|^2 \le \psi(u) \quad \forall u \in U,


$$


где $u^* \in U^*$.


\end{thms}





{\bf Доказательство.} 


Выберем произвольно число $\mu \in (0,1]$ и точку $u^* \in U^*$


и определим\linebreak


$v_\mu=\mu u^* + (1-\mu)u$.


Учитывая выпуклость функций $h(u, \cdot)$ и $f$, имеем


\begin{multline*}


\psi(u) \ge -h(u, v_\mu) + f(u) - f(v_\mu) - 0.5\alpha\|u-v_\mu\|^2 


\ge -\mu h(u, u^*) -(1-\mu)h(u,u) + \\


+ f(u) - \mu f(u^*) - (1-\mu)f(u) - 0.5\alpha\mu^2\|u-u^*\|^2 = \\


= -\mu h(u, u^*) +\mu [f(u) - f(u^*)] - 0.5\alpha\mu^2\|u-u^*\|^2.


\end{multline*}


Далее, из выпуклости функции $h(u,\cdot)$, 


свойства $h(u,u)=0$ и соотношения (\ref{n1a}) следует


$$


-h(u, u^*) + f(u) -f(u^*) \ge


\langle \nabla_v h(u,u^*) - \nabla_v h(u^*,u^*), u-u^*\rangle.  


$$


Поэтому с учетом свойства (A3) получим


\begin{multline*}


\psi(u) \ge \mu\langle \nabla_v h(u,u^*) - \nabla_v h(u^*,u^*), u-u^*\rangle  


-0.5\alpha\mu^2\|u-u^*\|^2 \ge \\


\ge (\mu\tau - 0.5 \alpha\mu^2)\|u-u^*\|^2 \ge \sigma\|u-u^*\|^2,


\end{multline*}


где


$$


\sigma  =


\begin{cases}


\tau-0.5\alpha & \text{при } \tau \ge \alpha;\\


0.5\tau^2/\alpha & \text{при } \tau < \alpha.\ \ \square


\end{cases}


$$





Введем обозначение


$$


S(u)=\{ v \in U^n \mid \psi(v) \le \psi(u)\} \quad \forall u \in U^n.


$$





\begin{cors}\label{cor1}


Пусть выполняется предположение (A3). Тогда лебегово


множество $S(u)$ для любого $u \in U^n $ является ограниченным.


\end{cors}





\begin{cors}\label{cor2}


Пусть выполняется предположение (A3). Тогда задача \ref{n1}


имеет единственное решение.


\end{cors}





\begin{pf}


Поскольку функция $\psi$ непрерывна, то она достигает cвоего минимального 


значения на непустом компактном множестве $S(u')$, где $u'$ --- 


произвольная точка из $U$, т.\,е. существует решение задачи (\ref{minpsi}), 


которое, очевидно, является стационарной точкой. Из теоремы \ref{tst} 


следует, что существует решение задачи \ref{n1}. В силу теоремы 


\ref{terr} и предложения \ref{p14} это решение единственно.


\end{pf}





\begin{thms}[{\series{m}\shape{n}\selectfont глобальная сходимость}]


\label{thms2} 


Пусть выполняются предположения {\em (A1), (A2)} и {\em (A3)}. Тогда 


итерационная последовательность $\{u^k\}$, построенная методом спуска, 


сходится к единственному решению задачи $\ref{n1}$.


\end{thms} 





\begin{pf}


Для доказательства сходимости метода воспользуемся теоремой сходимости A из 


(\cite{Zan69}, гл.\,4, раздел~4.5). 


Прежде всего покажем, что выполняются 


все условия этой теоремы.





1) Все точки последовательности $\{u^k\}$ принадлежат компактному 


подмножеству $U$. Условие следует из следствия \ref{cor1}.





2) $\psi(u^{k+1}) < \psi(u^k)$ для $u^k \notin U^*$ и $u^k \in U^*$ в случае 


$\psi(u^{k+1}) \ge \psi(u^k)$. Условие следует непосредственно из 


теоремы \ref{tst}, леммы \ref{ldesc} и правила выбора шага в методе.





3) Алгоритмическое отображение $u \mapsto w(u)$ 


замкнуто в любой точке $u \notin U^*$.


Условие следует из непрерывности отображения $u \mapsto v(u)$,


непрерывности функции $\psi$ на $U$ и компактности множества $U \cap S(u^0)$.





Таким образом, все условия теоремы сходимости А 


(\cite{Zan69}, гл.\,4, раздел~4.5) выполняются, поэтому 


любая предельная точка последовательности $\{u^k\}$ является стационарной 


(по свойству~1) хотя бы одна такая точка существует). 


Следовательно, в силу теоремы \ref{terr} и следствия \ref{cor2}


последовательность $\{u^k\}$


сходится к единственному решению задачи \ref{n1}.


\end{pf}





\begin{thebibliography}{9}





\bibitem{BC84}


Байокки~К., Капело~А. {\it Вариационные и квазивариационные 


неравенства. Приложения к задачам со свободной границей}. 


-- М.: Наука, 1988. -- 448 с.





\bibitem{BO94}


Blum~E., Oettli~W. {\it From optimization  and variational


inequalities to equilibrium problems} // The Mathem. Student.


-- 1994. -- V.\,63 -- \No\,1. -- P.\,123--145.





\bibitem{BS96}


Bianchi~M., Schaible~S. 


{\it Generalized monotone bifunctions and equilibrium problems} //


J. Optim. Theory Appl. -- 1996. -- V.\,90. -- \No\,1. -- P.\,31--43.





\bibitem{KS00}


Konnov~I.V., Schaible~S.


{\it Duality for equilibrium problems under generalized monotonicity} //


J. Optim. Theory Appl. -- 2000. -- V.\,104. -- \No\,2. -- P.\,395--408.





\bibitem{Pat97}


Patriksson~M. {\it  Merit functions and descent algorithms for a class


of variational inequality problems} //


Optimization. -- 1997. -- V.\,41. -- \No\,1. -- P.\,37--55.





\bibitem{DR90}


Демьянов В.Ф., Рубинов А.И. {\it Основы негладкого анализа и 


квазидифференциальное 


исчисление}. -- М.: Наука, 1990. -- 432~с.





\bibitem{Zan69}


Зангвилл У. {\it Нелинейное программирование. Единый подход.} 


-- М.: Советское радио, 1973. -- 312~с.





\end{thebibliography}





\vskip 2cm





\post{\it Казанский государственный}{\it Поступила}


\post{\it университет}{15.04.2003}





\label{end}


\end{document}


